\subsection{镜像对称}

当 $X$ 是 Calabi--Yau 流形时，\footnote{这一整套故事可以推广到 K\"ahler 流形，但为简单起见，我们不这样做。出于同样的理由，我们也不讨论超势（superpotential）在二维与三维镜像对称中所起的作用。}存在一个物理上的构造，给出以 $X$ 为目标空间的二维 $\mathcal N=(2,2)$ 超对称 sigma 模型 $\mathcal{T}^{d=2}(X)$。若关于 $Q_{\operatorname{A/B}}$ 作扭转，所得的拓扑量子场论（topological quantum field theory，TQFT）就称为 A/B 模型 $\mathcal{T}^{d=2}_{\operatorname{A/B}}(X)$。A 模型依赖于 $X$ 的辛拓扑，而 B 模型依赖于 $X$ 的复几何。

在这类超对称 sigma 模型的集合上，存在一个非凡的对偶，称为\textbf{镜像对称}（mirror symmetry）；它把 $\mathcal{T}^{d=2}(X)$ 与 $\mathcal{T}^{d=2}(X^\vee)$ 等同起来，其中 $X^\vee$ 是另一个 Calabi--Yau 流形，称为 $X$ 的\textbf{镜像对偶}（mirror dual）。

此外，这个对偶与拓扑扭转相容：$\mathcal{T}^{d=2}(X)$ 与 $\mathcal{T}^{d=2}(X^\vee)$ 之间的等同，与 $d=2$、$\mathcal N=(2,2)$ 超对称代数的一个对合相容，该对合交换 $Q_{\operatorname{A}}$ 与 $Q_{\operatorname{B}}$。因此，两个 $d=2$ 的 TQFT $\mathcal{T}^{d=2}_{\operatorname{A}}(X)$ 与 $\mathcal{T}^{d=2}_{\operatorname{B}}(X^\vee)$ 应当等价。这一想法带来了若干惊人的预言。其中最著名的是：五次三维流形（quintic threefold）上次数为 $d$ 的有理曲线的条数，若理解为 $\mathcal{T}^{d=2}_{\operatorname{A}}(X)$ 的关联函数（correlation function），应当等于 $\mathcal{T}^{d=2}_{\operatorname{B}}(X^\vee)$ 的关联函数，而后者更容易计算。

二维镜像对称（在数学中径称「镜像对称」）的非凡成功，促使人们为 $d=3$、$\mathcal N=4$ 的超对称场论考虑一个类似的对偶，称为\textbf{三维镜像对称}（3-dimensional mirror symmetry）：它把两个表面上很不相同的理论等同起来，记作 $\mathcal{T}$ 与 $\mathcal{T}^\vee$。与先前一样，超 Poincar\'e 代数中仍有两个有趣的拓扑扭转 $Q_{\operatorname{A}}$ 与 $Q_{\operatorname{B}}$，并且 $\mathfrak{a}$ 有一个交换二者的自同构。按照同样的逻辑，我们得到 $\mathcal{T}_{\operatorname{A}}$ 与 $\mathcal{T}_{\operatorname{B}}^{\vee}$ 作为拓扑扭转理论的一个等价；为了强调 TQFT 的观点，我们分别把它记作 $Z_{\operatorname{A}}^{\mathcal{T}}$ 与 $Z_{\operatorname{B}}^{\mathcal{T}^\vee}$。这看似纯粹抽象，但正如 $d=2$ 的情形，我们可以从这些拓扑扭转理论中得到更熟悉的数学对象。

我们将着重理解代数 $\mathcal{A}_{\operatorname{A/B}}(\mathcal{T})=Z_{\operatorname{A/B}}^{\mathcal{T}}(S^2)$。如第 \ref{sec:TQFT} 节所述，代数簇
\[\mathfrak{M}_{\operatorname{A}}(\mathcal{T})=\Spec Z_{\operatorname{A}}^{\mathcal{T}}(S^{2})\qquad \mathfrak{M}_{\operatorname{B}}(\mathcal{T})=\Spec Z_{\operatorname{B}}^{\mathcal{T}}(S^{2})\]
分别是相应理论的真空模空间（moduli space of vacua）。我们把它们称为理论 $\mathcal{T}$ 的\textbf{库仑分支} $\mathfrak{M}_{\operatorname{A}}(\mathcal{T})$ 与\textbf{希格斯分支} $\mathfrak{M}_{\operatorname{B}}(\mathcal{T})$；在物理文献中，这些对象通常有稍微不同的描述，但对我们所关心的理论而言，两者是一回事。

当然，把一个理论中的局部算子 $\mathcal{A}_{\operatorname{A}}(\mathcal{T})$ 与镜像理论中的 $\mathcal{A}_{\operatorname{B}}(\mathcal{T}^{\vee})$ 等同起来，在此情形下同样是镜像对称最重要的特征之一：
\[\mathfrak{M}_{\operatorname{A}}(\mathcal{T})\cong \mathfrak{M}_{\operatorname{B}}(\mathcal{T}^{\vee})\qquad \qquad \mathfrak{M}_{\operatorname{B}}(\mathcal{T})\cong \mathfrak{M}_{\operatorname{A}}(\mathcal{T}^{\vee}).\]
因此，我们称簇 $\mathfrak{M}_{\operatorname{A}}(\mathcal{T})$ 与 $\mathfrak{M}_{\operatorname{B}}(\mathcal{T})$ 互为\textbf{镜像}；按 \cite{BLPWgco} 的术语，它们互为辛对偶（symplectic dual）。这些簇的长处在于，它们属于人们熟悉的数学对象类型，同时又仍承载着理论 $\mathcal{T}$ 的许多结构。

% [原文 p11–p12]
